\subsection{Cross-Design Test: Real Covariates in the ABM}\label{sec:robustness-crossdesign}

Section~\ref{sec:method-estimators} flags an unresolved confound: the two estimators differ at once in modeling paradigm (household choice versus loan-level survival) and in data source (synthetic population versus real Freddie Mac loan records), so a low ABM share against a high hazard-framework share is equally consistent with either. To make partial progress on separating them, I re-ran the ABM's decision rule on structural covariates --- coupon, loan age, original LTV --- from Section~\ref{sec:hazard}'s 75,000-loan Freddie Mac sample, replacing the log-normal, FRED-median-calibrated draws used elsewhere. Behavioral draws (mobility desire, loss aversion, patience, transaction cost) stay synthetic, having no analogue in that data; the test therefore varies real structural heterogeneity --- and, until the post-stratification below, book composition jointly --- not a fully real population.

The ABM's mobility-scale parameter is calibrated to hold a real-world baseline turnover floor fixed (Appendix~\ref{sec:method-abm}), so a population with different LTV and coupon moments raises a question the covariate swap does not resolve on its own: re-derive the turnover-floor anchor on the new moments, or hold it fixed? I report both. The recalibrated variant applies Appendix~\ref{sec:method-frictions}'s recalibration rule, re-anchoring the mobility scale to reproduce the same 4--5\% turnover floor under the real population's moments; the frozen variant suspends that rule and leaves the original synthetic-population calibration unchanged.

I fixed the interpretive thresholds ex ante. A recovery above 50\% of the empirical benchmark would undercut the paradigm interpretation advanced in Sections~\ref{sec:abm}--\ref{sec:conclusion}; a recovery below 35\% would corroborate it, consistent with the synthetic-population ABM's own specifications; between 35\% and 50\% would be ambiguous.

The recalibrated variant recovers \$453.5 billion, 59.3\% of the benchmark, above that threshold; the frozen variant recovers \$159.5 billion, 20.9\%, inside the corroborating band (Table~\ref{tab:crossdesign} and Figure~\ref{fig:crossdesign}, both against the synthetic-population control). Both are single draws. A pre-committed fifty-seed sweep over both randomization layers, behavioral draws and covariate subsample (\texttt{cross\_design\_seeds}), shows the crossing is the estimator's, not the seed's: the recalibrated variant averages 60.2\%, standard deviation 3.3 points, 95\% band 55.1--66.0\%, the committed draw at the 48th percentile of its own distribution, all fifty seeds classifying as undercutting. The frozen variant averages 22.0\%, a band wide enough to report --- 10.9 points against the 7.7-point disclosure threshold fixed before the run --- but too narrow to reach the calibration swing. The operative uncertainty is calibration choice, not seed noise. The sweep also separates the two on the deep-gap floor: recalibrated CPR at an 8\% market rate averages 4.59\%, inside the observed 4--5\% turnover band on all fifty seeds; frozen averages 7.05\%, outside it on all fifty (range 6.31--7.63\%). Applied consistently, the admissibility test I use on the externally anchored variant below counts against the frozen leg: the leg reported as within the corroborating band is the one that misses the turnover moment. I read its 20.9\% as the less admissible of the two calibrations, not as corroboration, which strengthens the undercutting reading. And ``frozen'' is convention, not data: re-deriving that anchor per seed gives mean 43{,}493, standard deviation 1{,}181, minimum 39{,}984 against the committed 43{,}883, the very top of its own range.

\begin{figure}[H]
\centering
\includegraphics[width=\textwidth]{fig5_cross_design}
\caption{The cross-design test as paired panels. Panel (a): dollar recovery for the two real-covariate variants against the pre-committed thresholds and the synthetic-population control band; the recalibrated variant's 59.3\% crosses the undercutting threshold, the frozen variant's 20.9\% stays within the corroborating band. Panel (b): lag-0 path correlations, all negative --- neither variant reproduces the empirical path's timing regardless of calibration choice, and after Section~\ref{sec:pathb}'s detrending result no estimator in this paper does.}
\label{fig:crossdesign}
\end{figure}

\begin{table}[H]
\centering
\small
\begin{threeparttable}
\caption{Cross-design test results: real Freddie Mac structural covariates substituted into the ABM under two calibration choices, against the synthetic-population production ABM.}
\label{tab:crossdesign}
\begin{tabular}{@{}p{4.0cm}lll@{}}
\toprule
 & Recalibrated (primary) & Frozen (robustness) & Synthetic control \\
\midrule
Mobility-scale parameter & 36,086 & 43,883 & 43,883 \\
Trapped liquidity & \$453.5B (59.3\%; fifty-seed mean 60.2\%) & \$159.5B (20.9\%; 22.0\%) & \$91.0B (11.9\%) \\
Mean CPR & 8.14\% & 11.63\% & 11.68\% \\
CPR path corr. (lag 0) & $-0.336$ & $-0.311$ & $-0.318$ \\
Peak lag & 0 (no positive alignment) & 0 (no positive alignment) & 0 \\
\bottomrule
\end{tabular}
\begin{tablenotes}[flushleft]\footnotesize
\item Peak lag follows the max-$|r|$ convention of Table~\ref{tab:estimators}. Under the alternative most-positive-$r$ convention, the synthetic control's best alignment is $+0.195$ at lag $-3$, the ``$+0.19$ at the three-month lead'' of Section~\ref{sec:abm-results}; the two conventions summarize the same cross-correlation profile, and the substantive conclusion of no positive contemporaneous alignment holds under either.
\end{tablenotes}
\end{threeparttable}
\end{table}

By my own ex-ante criterion the primary recalibrated variant moves against the paradigm interpretation: with real structural covariates and a like-for-like turnover anchor, the household-choice ABM explains far more of the benchmark than the synthetic-population specifications of Sections~\ref{sec:abm} and~\ref{sec:robustness-scale} suggested. But the \$294 billion gap between the two variants is produced solely by how the turnover-floor anchor is set, covariates fixed: the ABM's aggregate share is highly sensitive to a calibration choice orthogonal to data source, so the test cannot on its own cleanly attribute Sections~\ref{sec:abm} and~\ref{sec:hazard}'s contrast to modeling paradigm rather than to data source --- recalibration choice does comparable work to whichever the test was designed to isolate.

I therefore withdraw the paradigm reading as a claim of this paper. The threshold was pre-committed before the run, and it was crossed. A reading that survives only by narrowing, after the fact, what the threshold was taken to test is not one I will defend --- that narrowing is what the pre-commitment architecture exists to prevent. Two results the ABM sections do establish stand without it, neither a claim about modeling paradigm: no synthetic-population household-choice specification tested here explains the Federal Reserve's cap shortfall, and no estimator in this paper reproduces the monthly path once the window's shared trend is removed.

A pre-committed post-stratification separates composition from heterogeneity. I reweight the real-covariate population to the SOMA coupon-by-vintage-group composition, the book the benchmark is measured on, and re-derive the mobility scale under the same floor rule. That moves the recalibrated variant from 59.3\% to 76.3\% and the frozen from 20.9\% to 12.6\%; the reweight uses importance weights, effective sample size 1{,}332 of 10{,}000, disclosed. A seed-42 re-draw within SOMA-weighted cells gives 70.6\% recalibrated, 36.0\% frozen. Composition mismatch therefore understated rather than manufactured the recalibrated recovery: at book composition, real-covariate household choice with a like-for-like floor anchor explains roughly seventy to seventy-six percent of the benchmark. The ex-ante undercutting verdict survives composition adjustment, and the calibration-choice spread widens to roughly \$488 billion. The reweight is a level exercise.

A third variant drops the turnover-floor anchor entirely and takes both author-calibrated behavioral elements from Path B's external elasticity source: the mobility-desire scale from \citepos{liebersohn2024} zero-gap moving level, a 1.5\% quarterly ZIP-code moving hazard at an essentially zero average rate gap, 5.9\% annualized; the wait-and-see freeze share from what their 6.5\%-per-quarter-per-100 bp mobility decline implies at the production 150 bp trigger, 18.3\% against the author-calibrated 20\%. The DTI wall and loss-aversion coefficient are already externally sourced. Recovery does rise without the floor recalibration, to 150.6\% of benchmark, but only by predicting essentially zero prepayment at QT-window gaps: CPR at an 8\% market rate is 0.0\%, against observed 4--5\% discount-cohort turnover and the roughly 1.44--1.45\% quarterly moving hazard \citet{liebersohn2024} themselves observe at 2022--2023 gaps. The externally anchored gate structure can match the zero-gap moving level or the observed deep-gap turnover, not both; the exponential mobility-desire distribution that fits the former (scale 4{,}472 against the recalibrated 36{,}086) shuts moving down once the par-payoff penalty binds. I read that as clarifying the floor anchor's role: it keeps the ABM on observed turnover data rather than suppressing its rate response. The anchored dimension is never claimed as fit, since matching the floor is tautological by construction (Appendix~\ref{sec:method-frictions}); what the variant measures is what that tautology purchases, empirical admissibility at deep gaps, another instance of this subsection's deflate-as-well-as-inflate indeterminacy (same harness).

Scored against the pre-committed bands, as its own specification required, its 150.6\% classifies as undercutting the paradigm reading, and the frozen artifact records that classification. I do not read it as informative about the paradigm question, and I label that reading post hoc: the pre-run specification committed only to reporting the deep-gap floor diagnostic ``alongside whatever the recovery number is, whichever direction it moves,'' not to removing the recovery from the band when the diagnostic fails. The admissibility precondition---an empirically admissible deep-gap turnover level before a recovery is interpretable---was formulated after the result was seen, and I report it as a labelled post-run qualification, not a pre-committed exclusion. A reader who declines it should score the variant at 150.6\% against the 50\% band, which strengthens the undercutting reading this subsection concedes; no headline quantity turns on the choice.

I ran a minimal simulated-minimum-distance calibration to reconcile the two external moments: a joint two-parameter fit of the mobility-desire scale plus a forced-move point mass added to the structural family, with moments, admissibility bands, distance and verdict vocabulary fixed before execution and a parity gate replaying the committed 59.303\% first. The verdict is joint infeasibility, now with two degrees of freedom, not one: at the distance-minimizing point (scale 1{,}000, mass 3.9\%) the zero-gap moving level overshoots its band, 0.0644 against $[0.055, 0.063]$, just as the deep-gap turnover reaches its own, 0.0394 against $[0.040, 0.050]$, the forced mass lifting both together at the scale's floor. The near-fit recovers 107.1\% of benchmark through the same harness, far above the pre-committed cross-design bands: external anchoring pushes this family past the benchmark, not toward it. I deliberately do not fit QT-window CPR moments; the recovery statistic is a functional of that path, so fitting it would make the test a fitted quantity---the overfitting the pre-commitment architecture exists to prevent. Hazard-side headline quantities are not functions of ABM calibration and cannot move under any outcome of this run.

Both variants retain the ABM's central path-level failure (Table~\ref{tab:crossdesign}); Path B's distinction is level accuracy, not timing. The timing critique of Section~\ref{sec:abm} therefore survives whichever calibration variant is preferred. The dollar-recovery critique does not straightforwardly survive under the recalibrated variant.
